\documentclass[10pt,a4paper]{amsart}
\usepackage[margin=19mm]{geometry}
\usepackage{amsmath,amssymb,mathtools}
\usepackage[scr=rsfs,cal=euler]{mathalfa}
\usepackage{xfrac}
\usepackage[unicode,hidelinks,bookmarks=false]{hyperref}
\newcommand{\ie}{i.e.\ }
\newcommand{\cf}{cf.\ }
\newcommand{\resp}{respectively\ }
\newcommand{\Subsection}{\subsection}
\providecommand{\nobreakdash}{\nobreak\hbox{-}}
\setlength{\emergencystretch}{2em}
\multlinegap0pt
\usepackage{amssymb}
\usepackage{xcolor}
\usepackage{mathtools}
\usepackage{tikz}
\usepackage{algorithm}\usepackage{algpseudocode}
\usepackage{enumerate}
\newtheorem{theorem}{Theorem}
\title{\bf Structure-Preserving Learning of Nonholonomic Dynamics}
\author{Thomas Beckers, Anthony Bloch, Leonardo Colombo}
\date{Source: arXiv:2603.27580v1 (CC BY 4.0)}
\begin{document}
\maketitle

\noindent\emph{Source and licence.} \bf Structure-Preserving Learning of Nonholonomic Dynamics. Version arXiv:2603.27580v1 (CC BY 4.0), distributed by arXiv under the Creative Commons Attribution 4.0 International licence, http://creativecommons.org/licenses/by/4.0/, as recorded in the arXiv licence metadata for this exact version and shown on the abstract page https://arxiv.org/abs/2603.27580v1. Full licence notice: \texttt{licenses/arxiv-2603.27580-CC-BY-4.0.txt}.\par\smallskip

\noindent\textit{Editorial scope.} Counted unit: the article's theorem theorem (source0.tex lines 588--595). The article's complete original proof of it is delivered with it (source0.tex lines 597--657, one \texttt{proof} environment of 61 lines). No mathematics is written, repaired or re-derived: the statement and the proof are the article's own lines, in the article's own notation, and no retained line is substituted or re-flowed.  No further source-local context is needed: every cross-reference of every retained line resolves inside the two delivered spans.  Nothing was omitted from any retained span, so the package carries no omission record.  No retained line contains a citation, so the wrapper carries no bibliography and no bibliography entry.  No retained line contains a figure environment, an image-inclusion command or drawing code, so no image is delivered and the package declares no asset dependency.  Editorial formatting only, in addition: an AMS \texttt{amsart} preamble that restates the article's own declarations and macros used by the retained lines, namely the environments \texttt{theorem} on separate counters, together with the standard packages those lines need.  The retained environments print in this document as Theorem 1 in order of appearance, whereas the article prints the same items with the numbering of its own full document.  The complete native source of the article as distributed in the arXiv e-print for this exact version is bundled unchanged as \texttt{sources/197359/source0.tex}.
\begin{theorem}
Let $\hat f_N$ be the Gaussian process estimator obtained using the nonholonomic kernel $K_{\mathrm{NH}}$ and $N$ distinct training samples. Under the above assumptions, $\hat f_N$ is uniformly consistent in probability, namely, for every $\varepsilon>0$,
\[
\mathbb{P}\!\left(
\sup_{q\in Q}\|\hat f_N(q)-f^\star(q)\|>\varepsilon
\right)\to 0\,\, \text{as } N\to\infty.
\]
\end{theorem}

\begin{proof}
By Proposition 1, the RKHS associated with the nonholonomic kernel is
\[
\mathcal H_{\mathrm{NH}}
=
\left\{
f:Q\to\mathbb R^n
\;\middle|\;
f(q)=P(q)g(q),\ \ g\in\mathcal H_{kI_n}
\right\}.
\]
Since $f^\star\in\mathcal H_{\mathrm{NH}}$, there exists $g^\star\in\mathcal H_{kI_n}$ such that $f^\star(q)=P(q)g^\star(q)$,\,$q\in Q$.


Let $\hat g_N$ denote the Gaussian process estimator associated with the kernel
$K_0(q,q')=k(q,q')I_n$. Under Assumptions 1 and 2 and standard consistency results for kernel-based GP regression, $\hat g_N$ is uniformly consistent in probability for $g^\star$, namely, for every $\varepsilon>0$,
\[
\mathbb{P}\!\left(
\sup_{q\in Q}\|\hat g_N(q)-g^\star(q)\|>\varepsilon
\right)\to 0
\quad\text{as }N\to\infty.
\]

On the other hand, the estimator induced by the nonholonomic kernel satisfies
$\hat f_N(q)=P(q)\hat g_N(q)$. Therefore, for every $q\in Q$, $\hat f_N(q)-f^\star(q)
=
P(q)\bigl(\hat g_N(q)-g^\star(q)\bigr)$,
and hence $\|\hat f_N(q)-f^\star(q)\|
\le
\|P(q)\|\,\|\hat g_N(q)-g^\star(q)\|$.


By construction, $P(q)=I-A(q)^\dagger A(q)$ is the orthogonal projector onto  $\mathcal D_q=\ker A(q)$. In particular, $P(q)$ is self-adjoint and idempotent, that is, $P(q)^\top=P(q)$ and $P(q)^2=P(q)$. Hence, with respect to the Euclidean norm, its operator norm satisfies $\|P(q)\|\le 1$, $q\in Q$. 

Indeed, if $v\in\mathbb R^n$ is decomposed orthogonally as $v=v_{\mathcal D}+v_{\mathcal D^\perp}$,
$v_{\mathcal D}\in\mathcal D_q,\, v_{\mathcal D^\perp}\in\mathcal D_q^\perp$,
then $P(q)v=v_{\mathcal D}$, and therefore
$\|P(q)v\|=\|v_{\mathcal D}\|\le \|v\|.$ Consequently,
\[
\sup_{q\in Q}\|\hat f_N(q)-f^\star(q)\|
\le
\sup_{q\in Q}\|\hat g_N(q)-g^\star(q)\|.
\]

Now let $\varepsilon>0$. Then
\[
\mathbb{P}\!\left(
\sup_{q\in Q}\|\hat f_N(q)\!-\!f^\star(q)\|\!>\!\varepsilon\!
\right)
\!\le\!
\mathbb{P}\!\left(
\sup_{q\in Q}\|\hat g_N(q)\!-\!g^\star(q)\|\!>\!\varepsilon\!
\right)\!.
\]
Since the right-hand side converges to zero as $N\to\infty$, we conclude that, for every $\varepsilon>0$,
\[
\mathbb{P}\!\left(
\sup_{q\in Q}\|\hat f_N(q)-f^\star(q)\|\>\varepsilon
\right)\to 0
\quad\text{as }N\to\infty.
\] \end{proof}

\end{document}
